 >>?% ======================================================================
% CAPITULO 1 --- Introducao: Por que Arbitragem?
% ======================================================================
\chapter{Introducao: Por Que Arbitragem?}

\section{O Conceito Mais Simples do Mundo}

Imagine que voce esta no aeroporto de Guarulhos. Na casa de cambio A, 1 dolar
americano custa R\$ 5,00. Na casa de cambio B, a 20 metros de distancia, 1
dolar custa R\$ 4,90. O que voce faz?

\begin{enumerate}
  \item Compra 1 dolar na casa B, pagando R\$ 4,90.
  \item Imediatamente vende este mesmo dolar na casa A, recebendo R\$ 5,00.
  \item Lucro liquido: R\$ 0,10.
\end{enumerate}

\textbf{Risco: zero.} As duas transacoes sao simultaneas. Voce nao precisa
de capital inicial --- pode pedir emprestado os R\$ 4,90 e devolver
instantaneamente com os R\$ 5,00 recebidos. \textbf{Investimento inicial:
zero.}

Este e o exemplo mais simples de \textbf{arbitragem}: lucro sem risco obtido
pela exploracao de uma inconsistencia de precos. A palavra vem do frances
\textit{arbitrer} (julgar, decidir) e do latim \textit{arbitrari} (dar um
veredito). O arbitrador e aquele que ``julga'' que dois precos deveriam ser
iguais e age sobre essa discrepancia.

\section{Arbitragem no Mundo Real}

Em mercados financeiros reais, oportunidades como a do exemplo duram
milissegundos --- ou menos. Algoritmos de \textit{high-frequency trading}
(HFT)\footnote{Aldridge, I. \textit{High-Frequency Trading: A Practical Guide
to Algorithmic Strategies and Trading Systems}. Wiley, 2nd ed., 2013. ISBN:
978-1118343500. \textbf{Resenha:} Referencia pratica sobre estrategias de HFT,
incluindo arbitragem estatistica, market making e execucao algoritmica.} sao
programados para detectar e executar estas oportunidades antes que qualquer
humano possa reagir. Em 2023, estimava-se que 50-60\% do volume de negociacao
nas bolsas americanas era gerado por algoritmos de HFT.

Mas a \textit{teoria} da arbitragem --- e nao sua pratica --- e o alicerce sobre o
qual todo o edificio das financas quantitativas modernas foi construido. A
hipotese de \textbf{ausencia de arbitragem} e o ponto de partida para:

\begin{itemize}
  \item O modelo de Black-Scholes-Merton (1973) para aprecamento de opcoes
  \item A teoria de carteiras de Markowitz (1952)
  \item O Capital Asset Pricing Model (CAPM)
  \item Os modelos de estrutura a termo de taxas de juros
  \item Os modelos de risco de credito
  \item Praticamente toda a matematica financeira moderna
\end{itemize}

\section{Por Que a Ausencia de Arbitragem e uma Hipotese Razoavel?}

A resposta e surpreendentemente simples: \textbf{se houvesse uma oportunidade
de arbitragem obvia e sem risco, alguem ja a teria explorado ate ela
desaparecer.}

Este argumento, conhecido como \textbf{principio da nao arbitragem}, e um
caso especial da \textbf{hipotese do mercado eficiente}
(HME)\footnote{Fama, E.F. \textit{Efficient Capital Markets: A Review of Theory
and Empirical Work}. Journal of Finance, 25(2):383--417, 1970. DOI:
10.2307/2325486. \textbf{Resenha:} O artigo seminal que formalizou a HME em
tres formas: fraca (precos refletem informacao historica), semiforte (precos
refletem toda informacao publica) e forte (precos refletem toda informacao,
inclusive privada). Fama recebeu o Premio Nobel de Economia em 2013.}: em um
mercado eficiente, os precos refletem toda a informacao disponivel. Se houvesse
uma discrepancia que permitisse lucro sem risco, os proprios agentes que a
exploram fariam os precos convergirem --- eliminando a oportunidade.

E importante notar que a HME \textbf{nao} afirma que os mercados sao perfeitos.
Ela afirma que e \textbf{dificil} bater o mercado consistentemente apos
descontar custos de transacao e risco. A arbitragem pura (risco zero, lucro
certo) e virtualmente inexistente em mercados liquidos. A \textbf{arbitragem
estatistica} (risco muito baixo, lucro muito provavel) e o que os fundos
quantitativos perseguem.

\section{Arbitragem, na Arbitragem, e Quase-Arbitragem}

Precisamos distinguir tres conceitos que sao frequentemente confundidos:

\begin{defi}[Arbitragem Pura]
Uma estrategia autofinanciada que:
\begin{enumerate}[(i)]
  \item Nao requer investimento inicial (ou requer investimento negativo ---
      voce recebe dinheiro para entrar na estrategia),
  \item Nunca gera perda (payoff $\geq 0$ com probabilidade 1),
  \item Tem probabilidade positiva de lucro ($P[\text{payoff} > 0] > 0$).
\end{enumerate}
\end{defi}

\begin{defi}[Quase-Arbitragem]
Uma estrategia que ``quase'' satisfaz as condicoes acima --- o risco nao e
exatamente zero, mas tende a zero assintoticamente. Estas estrategias sao
excluidas pela condicao NFLVR (No-Free-Lunch-with-Vanishing-Risk).
\end{defi}

\begin{defi}[Arbitragem Estatistica]
Uma estrategia cujo valor esperado e positivo e cuja probabilidade de perda
tende a zero quando o horizonte de tempo tende a infinito. Risco existe, mas
e controlado. E o dominio dos fundos quantitativos.
\end{defi}

\section{O Primeiro Teorema Fundamental --- Versao Intuitiva}

O \textbf{Primeiro Teorema Fundamental do Aprecamento de Ativos} (FTAP) e o
resultado central da teoria de arbitragem. Em sua versao mais intuitiva:

\begin{center}
\fbox{\begin{minipage}{0.88\textwidth}
\textbf{FTAP (versao intuitiva):} Em um mercado livre de arbitragem, existe
uma medida de probabilidade $P^*$ --- chamada \textbf{medida martingal
equivalente} ou \textbf{medida neutra ao risco} --- sob a qual o preco descontado
de todo ativo e um martingal:
\[
\hat{S}_t = \mathbb{E}^{P^*}_t[\hat{S}_T], \quad \forall T \geq t
\]
Em palavras: o melhor previsor do preco futuro (descontado) e o preco de hoje.
E, reciprocamente, se tal medida existe, nao pode haver arbitragem.
\end{minipage}}
\end{center}

\textbf{Traducao para o portugues claro:} Se eu consigo atribuir probabilidades
aos cenarios futuros de tal forma que, em media, nenhum ativo tende a subir
ou cair (todos sao ``jogos justos''), entao e impossivel montar uma estrategia
que lucre sem risco. E vice-versa: se nao ha arbitragem, entao \textit{existe}
uma tal atribuicao de probabilidades --- mesmo que ela nao corresponda as
probabilidades reais.

\section{Um Exemplo Concreto: Aprecamento de Opcao}

Suponha uma acao que hoje custa $S_0 = 100$. Em um mes, dois cenarios sao
possiveis:
\[
S_1 = \begin{cases}
120 & \text{com probabilidade real } p = 0,5 \\
80  & \text{com probabilidade real } 1-p = 0,5
\end{cases}
\]
A taxa de juros e zero (para simplificar).

Considere uma \textbf{opcao de compra} (\textit{call}) com preco de exer
$K = 105$ e vencimento em 1 mes. Seu payoff e:
\[
C_1 = \max(S_1 - K, 0) = \begin{cases} 15 & \text{se } S_1 = 120 \\
0 & \text{se } S_1 = 80 \end{cases}
\]

\textbf{Abordagem ingenua (errada):} O valor esperado do payoff e
\[
\mathbb{E}^P[C_1] = 0,5 \times 15 + 0,5 \times 0 = 7,50
\]
Logo o preco justo seria R\$ 7,50. Certo? \textbf{Errado.}

Se a opcao custasse R\$ 7,50, eu poderia montar uma arbitragem:
\begin{enumerate}
  \item Vendo 1 opcao por R\$ 7,50.
  \item Compro $x$ acoes por R\$ 100x.
  \item Ajusto $x$ para que o valor da carteira seja identico nos dois
        cenarios (eliminando o risco):
        \begin{align*}
          120x - 15 &= 80x - 0 \\
          40x &= 15 \\
          x &= 0,375
        \end{align*}
  \item Valor final da carteira em qualquer cenario:
        \[
        120 \times 0,375 - 15 = 45 - 15 = 30
        \]
        \[
        80 \times 0,375 - 0 = 30 - 0 = 30
        \]
  \item Custo da carteira hoje: $100 \times 0,375 - 7,50 = 37,50 - 7,50 = 30$.
  \item Valor em 1 mes: 30 (certo).
  \item Lucro: $30 - 30 = 0$? Nao --- a arbitragem esta em outro lugar.
\end{enumerate}

Na verdade, a verdadeira arbitragem e mais sutil: se a opcao custa R\$ 7,50,
o retorno da carteira coberta e $30/30 = 0\%$ --- a taxa livre de risco. Mas
ha uma carteira que custa menos e produz o mesmo payoff:

\begin{itemize}
  \item Compre $0,375$ acoes por R\$ 37,50.
  \item Tome emprestado R\$ 30,00 a taxa zero.
  \item Custo liquido: R\$ 7,50.
  \item Payoff: $45 - 30 = 15$ ou $30 - 30 = 0$ --- replica a opcao.
\end{itemize}

Para eliminar a arbitragem, o preco da opcao deve ser exatamente o custo da
carteira replicante. A medida martingal equivalente $P^*$ atribui
probabilidade $p^*$ ao cenario de alta tal que:

\[
100 = p^* \times 120 + (1-p^*) \times 80 \implies p^* = 0,5
\]

(Neste caso, com taxa de juros zero, $p^* = 0,5$.) O preco justo da opcao e:

\[
C_0 = \mathbb{E}^{P^*}[C_1] = 0,5 \times 15 + 0,5 \times 0 = 7,50
\]

Espere --- isso da 7,50 de novo! O que mudou? A \textbf{medida} mudou. Na medida
real $P$, o retorno esperado da acao e $0,5 \times 120 + 0,5 \times 80 = 100$,
que e exatamente o preco de hoje. Mas a medida $P^*$ \textbf{nao e a medida
real} --- e a medida sob a qual o retorno esperado da acao e a taxa de juros.
Com taxa zero, $P^* = P$ neste exemplo simplificado, mas em geral $P^* \neq P$.

\section{Black-Scholes-Merton: A Revolucao de 1973}

O exemplo acima e um modelo binomial de um periodo. A grande contribuicao de
Black, Scholes e Merton\footnote{Black, F. e Scholes, M. \textit{The Pricing of
Options and Corporate Liabilities}. Journal of Political Economy, 81(3):637--654,
1973. DOI: 10.1086/260062. \textbf{Resenha:} O artigo que revolucionou as
financas. Derivou a equacao diferencial parcial que governa o preco de opcoes
europeias. Black e Scholes usaram um argumento de hedge continuo (carteira
livre de risco) para derivar a EDP. Merton\footnotemark\ generalizou para
tempo continuo com processos de Ito. Scholes e Merton receberam o Premio Nobel
de Economia em 1997.}
\footnotetext{Merton, R.C. \textit{Theory of Rational Option Pricing}. Bell
Journal of Economics and Management Science, 4(1):141--183, 1973. DOI:
10.2307/3003143. \textbf{Resenha:} Estendeu Black-Scholes com tratamento
rigoroso em tempo continuo e provou que o preco de uma opcao e o valor esperado
descontado do payoff sob a medida martingal equivalente.} foi estender esta
ideia para tempo continuo com negociacao continua --- o que levou a famosa
\textbf{equacao de Black-Scholes}:

\[
\frac{\partial V}{\partial t} + \frac{1}{2}\sigma^2 S^2
\frac{\partial^2 V}{\partial S^2} + rS\frac{\partial V}{\partial S} - rV = 0
\]

A sol para uma opcao de compra europeia e a celebrada formula:

\[
C(S,t) = S\Phi(d_1) - Ke^{-r(T-t)}\Phi(d_2)
\]

onde $\Phi$ e a funcao de distribuicao normal padrao e
\[
d_{1,2} = \frac{\log(S/K) + (r \pm \sigma^2/2)(T-t)}{\sigma\sqrt{T-t}}
\]

\section{A Pergunta que a GAT Responde}

A teoria classica de arbitragem, que acabamos de resumir, e inteiramente
formulada na linguagem da \textbf{probabilidade}: medidas martingais
equivalentes, martingais, processos de Ito, equacoes diferenciais estocasticas.

A \textbf{Geometric Arbitrage Theory} (GAT), proposta por Ilinski
(2001)\footnote{Ilinski, K. \textit{Physics of Finance: Gauge Modelling in
Non-Equilibrium Pricing}. Wiley, 2001. ISBN: 978-0471877387. \textbf{Resenha:}
Primeiro livro a propor que arbitragem deveria ser modelada como curvatura de
uma conexao de gauge, em analogia com a eletrodinamica quantica. Ilinski ve
precos como campos de gauge e fluxos de capital como correntes. Embora menos
rigoroso que Farinelli, foi pioneiro na intuicao fisica para financas.} e
desenvolvida rigorosamente por Farinelli (2021)\footnote{Farinelli, S.
\textit{Geometric Arbitrage Theory and Market Dynamics Reloaded}. arXiv:
0910.1671v10, 2021.}, faz uma pergunta radicalmente diferente:

\begin{center}
\fbox{\begin{minipage}{0.85\textwidth}
\textbf{E se pudessemos \textit{ver} a arbitragem como uma propriedade
geometrica do mercado?} E se a ausencia de arbitragem fosse equivalente a
``planicidade'' de um espaco curvo --- como a ausencia de campo eletromagnetico
e equivalente a curvatura zero do fibrado $U(1)$?
\end{minipage}}
\end{center}

A resposta --- elegante e profunda --- e \textbf{sim}. E as implicacoes vao muito
alem da elegancia matematica:

\begin{itemize}
  \item \textbf{Classificacao de estrategias:} A holonomia do fibrado do
        mercado parametriza as classes de equivalencia de estrategias de
        arbitragem.
  \item \textbf{Hidrodinamica:} A condicao NFLVR e equivalente a uma equacao
        de continuidade --- a arbitragem flui como um fluido.
  \item \textbf{Leis de conservacao:} O Teorema de Noether estocastico conecta
        simetrias do mercado (invariancia sob mudanca de numerario) a
        quantidades conservadas (valor presente liquido esperado).
  \item \textbf{FTAP homotopico:} O Primeiro Teorema Fundamental ganha uma
        formulacao em termos de topologia diferencial --- grupos de homotopia e
        holonomia.
\end{itemize}

Nas proximas secoes, construiremos cada um destes conceitos a partir do zero.
Nao assumimos conhecimento previo de geometria diferencial. Comecaremos pela
probabilidade e pelo calculo estocastico --- os alicerces sobre os quais a GAT
se ergue.

\section{Estrutura Deste Modulo}

\begin{enumerate}[nosep]
  \item \textbf{Parte I --- Fundamentos (Caps. 1-4):} Revisao de probabilidade,
        calculo estocastico, e o modelo classico de mercado. O minimo
        necessario para entender a GAT.
  \item \textbf{Parte II --- GAT (Caps. 5-9):} Gauges, fibrados principais,
        conexoes, curvatura, e o Teorema Central (arbitragem = curvatura).
        O coracao da teoria.
  \item \textbf{Parte III --- Avancado (Caps. 10-13):} Hidrodinamica,
        formulacao Lagrangiana, topologia diferencial, e implementacoes
        computacionais. Para quem quer ir alem.
  \item \textbf{Apendices:} Revisao de geometria diferencial, derivadas de
        Nelson, exers resolvidos, e referencias completas.
\end{enumerate}




